% =====================================================================
\chapter{Proyecto de fase}
\label{cap:proyecto}
% =====================================================================

\section*{Informe: ¿siguen las palabras de un texto la ley de Zipf?}
\phantomsection\addcontentsline{toc}{section}{Informe: ¿siguen las palabras de un texto la ley de Zipf?}

\paragraph{Objetivo.} El estudiante redacta un informe matemático breve que combina demostraciones rigurosas con un pequeño estudio estadístico sobre un texto real. La pregunta es si la frecuencia $f_r$ de la palabra que ocupa el puesto $r$ en la lista de frecuencias se comporta como una ley de potencia $f_r\approx C\,r^{-s}$, con $s$ próximo a~$1$ (ley de Zipf~\cite{zipf}). El proyecto integra los nueve módulos de la fase y anticipa ideas del curso: las frecuencias de tokens (módulo~12), la entropía (módulo~6) y las leyes de escala ajustadas en escala log-log (módulo~16).

\paragraph{Datos.} Un texto en español de dominio público de al menos 20\,000 palabras, por ejemplo una novela del Proyecto Gutenberg~\cite{gutenberg}. Se pasa a minúsculas, se eliminan los signos de puntuación y se considera palabra toda secuencia maximal de letras. Debe indicarse la fuente exacta y el número de palabras ($n$) y de palabras distintas ($V$).

\paragraph{Parte teórica: demostraciones.} El informe debe incluir, redactadas con el estilo de los apuntes, las demostraciones siguientes.
\begin{enumerate}[label=D\arabic*.]
  \item \emph{Mínimos cuadrados.} Deducir los coeficientes $\widehat b_1=S_{xy}/S_{xx}$ y $\widehat b_0=\ybar-\widehat b_1\xbar$ anulando las derivadas parciales de $\SC(b_0,b_1)=\sum_i(y_i-b_0-b_1x_i)^{2}$ y demostrar que el punto crítico es el mínimo global cuando $S_{xx}>0$ (Módulo 0.7, §5.1 y §5.2; Módulo 0.9, §6.1).
  \item \emph{Media y varianza muestrales.} Demostrar que $\mathbb E[\Xbar]=\mu$, $\Var(\Xbar)=\sigma^{2}/n$ y $\mathbb E[S^{2}]=\sigma^{2}$ para una muestra aleatoria simple (Módulo 0.8, §5.2 y §5.3; Módulo 0.9, §2.4 y §2.5).
  \item \emph{Números armónicos.} Demostrar que $\ln(n+1)\le H_n\le1+\ln n$ para todo $n\ge1$, donde $H_n=\sum_{k=1}^{n}1/k$, comparando la suma con $\int\dd x/x$ (Módulo 0.6, §5.3). Deducir que, en el modelo $p_r=1/(r\,H_V)$, la probabilidad de las diez palabras más frecuentes es $H_{10}/H_V$ y acotarla con la desigualdad anterior.
  \item \emph{Linealización.} Demostrar que, para $C>0$, la relación $f=C\,r^{-s}$ entre números positivos equivale a $\ln f=\ln C-s\ln r$, y que la transformación $(r,f)\mapsto(\ln r,\ln f)$ es una biyección de $(0,\infty)^{2}$ en $\R^{2}$ (Módulo 0.2, §3.1; Módulo 0.3, §5.2).
  \item \emph{Chebyshev frente al intervalo normal.} Usar la desigualdad de Chebyshev para hallar un $n$ que garantice $\PP(\abs{\Xbar-\mu}\ge0{,}1)\le0{,}05$ en función de $\sigma$, y compararlo con el que da el intervalo normal (Módulo 0.8, §6.4; Módulo 0.9, §4.3).
\end{enumerate}

\paragraph{Parte práctica: estudio estadístico.}
\begin{enumerate}[label=E\arabic*.]
  \item \emph{Descripción.} Para la longitud en letras de las palabras: media, mediana, cuartiles, recorrido intercuartílico, desviación típica, histograma y diagrama de caja (Módulo 0.9, §1.2 a §1.4).
  \item \emph{Intervalo de confianza.} Intervalo del 95\,\% para la longitud media de las palabras, con una discusión de la hipótesis de independencia (Módulo 0.9, §4.4).
  \item \emph{Ajuste de Zipf.} Recta de mínimos cuadrados de $\ln f_r$ frente a $\ln r$ para $r=1,\dots,100$; estimación de $s$ y de $C$, coeficiente $R^{2}$ y gráfico de residuos. Discutir si $s$ es compatible con~$1$ y en qué rangos falla el ajuste (Módulo 0.9, §6.2 a §6.4).
  \item \emph{Contraste.} Contraste $\chi^{2}$ de bondad de ajuste de la hipótesis «las cinco vocales son equiprobables», con $\chi^{2}_{4;0{,}95}=9{,}488$ (Módulo 0.9, §5.5). Si se hacen varios contrastes, aplicar la corrección de Bonferroni (Módulo 0.9, §5.6).
  \item \emph{Predicción.} Comparar la proporción observada de apariciones de las diez palabras más frecuentes con la que predice el modelo ajustado y con la cota de~D3 (Módulo 0.8, §3.4).
\end{enumerate}

\paragraph{Entregables.}
\begin{enumerate}[label=\arabic*.]
  \item Informe de 8 a 12 páginas en \LaTeX, preferiblemente con la plantilla del curso, con las secciones: resumen, datos, demostraciones (D1 a D5), resultados (E1 a E5), discusión y limitaciones.
  \item Código (Python, R o una hoja de cálculo) que reproduzca todas las cifras y figuras a partir del texto original.
  \item Al menos tres figuras: histograma de longitudes, diagrama log-log de frecuencias con la recta ajustada y gráfico de residuos.
  \item Tabla resumen final con $n$, $V$, media y desviación típica de la longitud, intervalo de confianza, $\widehat s$, $\widehat C$, $R^{2}$, estadístico $\chi^{2}$ con su $p$-valor y la proporción observada y predicha del apartado~E5.
\end{enumerate}

\paragraph{Criterios de evaluación.} El proyecto se puntúa sobre 100 puntos con la rúbrica del \cref{tab:rubrica}. Se aprueba con 70 puntos o más.

\begin{table}[h]
\centering\footnotesize
\caption{Rúbrica del proyecto de fase.}\label{tab:rubrica}
\begin{tabularx}{\textwidth}{@{}>{\raggedright\arraybackslash}p{2.5cm}c>{\raggedright\arraybackslash}X>{\raggedright\arraybackslash}X>{\raggedright\arraybackslash}X@{}}
\toprule
Criterio & Peso & Excelente & Suficiente & Insuficiente \\
\midrule
Demostraciones D1--D5 & 35 & Las cinco completas, con cada paso justificado y las hipótesis explícitas & Cuatro correctas o cinco con lagunas menores & Menos de cuatro o con errores lógicos (por ejemplo, confundir una implicación con su recíproca) \\
\addlinespace
Estadística descriptiva (E1) & 10 & Todas las medidas y gráficos, bien interpretados & Medidas correctas con interpretación escasa & Medidas erróneas o sin gráficos \\
\addlinespace
Ajuste de Zipf (E3) & 20 & Ajuste correcto, $R^{2}$, residuos analizados y discusión de los rangos en que falla & Ajuste correcto sin diagnóstico de residuos & Ajuste en escala lineal o sin transformar los datos \\
\addlinespace
Inferencia y probabilidad (E2, E4, E5) & 20 & Intervalo, contraste y predicción correctos, con hipótesis discutidas & Cálculos correctos sin discutir las hipótesis & Grados de libertad, cuantiles o conclusiones erróneos \\
\addlinespace
Redacción y reproducibilidad & 15 & Informe claro, figuras rotuladas y código que reproduce todo & Informe legible y código parcial & Sin código o con cifras sin origen \\
\bottomrule
\end{tabularx}
\end{table}

\paragraph{Calificación del hito.} Nota del examen $=$ puntos$/10$; nota del proyecto $=$ puntos$/10$. Nota del hito $=0{,}6\cdot\text{examen}+0{,}4\cdot\text{proyecto}$. Para superar el hito hay que obtener al menos un 7 en cada una de las dos partes.
